\section{Epistemic transition envelopes}
\label{sec:escalation-theory}

The general object studied in this paper is not a particular diagnosis,
repair, rollback, or belief-revision component. It is an
\emph{epistemic transition envelope}: the admissible frontier of scientific
changes after those components have supplied their outputs. The envelope
separates three questions that are often collapsed: which transitions could
restore performance, which transitions the evidence authorizes, and which
authorized transition is minimal in the scientific change order. The
$K/W/M$ architecture used in the experiments is one realization of this
object, not part of its mathematical definition.

\subsection{Contexts, transition orders, and admissible fronts}

Let $\mathcal E$ be an arbitrary set of evidence contexts. A context $e$ has
a candidate transition space $\mathcal R_e$ and a strict partial order
$\prec_e$. The interpretation of $r'\prec_e r$ is that $r'$ makes strictly
less scientific change than $r$ in context $e$. The order need not be a
chain: revisions of a measurement model and a problem boundary, for example,
can be incomparable. We require only that the order restricted to any
admissible set encountered below be well founded. This includes every finite
revision space but also permits infinite spaces without infinite descending
chains.

Let $C_{1,e},\ldots,C_{k,e}$ be hard admissibility predicates. Their
conjunction, admissible set, and minimal front are
\[
 C_e(r)=\bigwedge_{i=1}^{k}C_{i,e}(r),\qquad
 \mathcal A(e)=\{r\in\mathcal R_e:C_e(r)\},\qquad
 \mathcal F(e)=\min_{\prec_e}\mathcal A(e).
\]
For the scientific-escalation instance evaluated here, $k=5$ and
\begin{align*}
J_e(r)&:\ \text{the evidence licenses revision at the layer changed by $r$},\\
T_e(r)&:\ \text{$r$ restores the same registered task under intervention},\\
P_e(r)&:\ \text{$r$ preserves every registered sibling or invariant},\\
D_e(r)&:\ \text{observed dependency impact equals the declared reopen scope},\\
B_e(r)&:\ \text{$r$ respects the registered action and resource budget}.
\end{align*}
Other scientific domains may replace or add hard predicates without changing
the results below.

Let $Q(e)$ say that the registered task has a material residual. The
transition decision $\Gamma$ is
\[
\Gamma(e)=
\begin{cases}
\mathsf{NO\_CHANGE}, & \neg Q(e),\\
\mathsf{UNRESOLVED}, & Q(e)\ \text{and}\ \mathcal A(e)=\varnothing,\\
r, & Q(e)\ \text{and}\ \mathcal F(e)=\{r\},\\
\mathsf{MEASURE}(\mathcal F(e)), & Q(e)\ \text{and}\
|\mathcal F(e)|>1.
\end{cases}
\]
The last output carries the unresolved front rather than hiding an arbitrary
tie-break. A later discriminating observation may refine the context and
order. Thus abstention and further measurement are identified scientific
actions, not missing values silently converted to failure or success.

\paragraph{Proposition 1 (sound, non-compensatory minimality).}
If $\Gamma(e)=r\in\mathcal R_e$, then every $C_{i,e}(r)$ holds and there is no
$r'\in\mathcal A(e)$ with $r'\prec_e r$. Moreover, optimizing any utility,
accuracy, or cost score over $\mathcal F(e)$ cannot select a transition that
fails a hard predicate, regardless of its score.

\paragraph{Proof.}
The first two statements are the definition of membership in
$\mathcal F(e)$. The last follows from
$\mathcal F(e)\subseteq\mathcal A(e)$: candidates outside the conjunction are
not in the optimization domain and cannot compensate for a failed gate.
\hfill$\square$

This is the minimal-escalation safety property tested by the mechanical
studies. The broader envelope formalism, however, is not tied to five gates,
three epistemic layers, or a finite generator.

\subsection{Exact substitution is a factorization property}

Suppose an arbitrary product of donor mechanisms exposes an interface
$\phi:\mathcal E\to\mathcal Z$. The interface can combine outputs from
diagnosis, repair, belief revision, dependency maintenance, model discovery,
planning, or human review. A downstream policy sees $\phi(e)$ but not $e$.
Write $\mathcal Y$ for the output space of $\Gamma$.

\paragraph{Theorem 1 (exact donor-substitution theorem).}
The following are equivalent:
\begin{enumerate}
  \item there exists a policy $g:\phi(\mathcal E)\to\mathcal Y$ such that
  $\Gamma=g\circ\phi$;
  \item $\Gamma$ is constant on every fibre of $\phi$, that is,
  \[
  \phi(e)=\phi(e')\quad\Longrightarrow\quad\Gamma(e)=\Gamma(e');
  \]
  \item the partition of contexts induced by $\phi$ refines the partition
  induced by the required transition decision.
\end{enumerate}

\paragraph{Proof.}
If $\Gamma=g\circ\phi$, equal interface values give equal decisions, so (1)
implies (2). If (2) holds, define $g(z)=\Gamma(e)$ for any $e$ with
$\phi(e)=z$. Fibre constancy makes this definition independent of the chosen
representative, proving (1). Statement (3) is the partition form of (2).
\hfill$\square$

The theorem gives a necessary and sufficient condition, not merely the
sufficient condition that a donor product reconstruct the full tuple
$(\mathcal R_e,\prec_e,C_{1,e},\ldots,C_{k,e},Q(e))$. Full reconstruction
guarantees equivalence, but a smaller statistic is enough whenever it separates
all contexts requiring different decisions. Conversely, component count,
mechanism sophistication, and successful repair rates cannot compensate for a
mixed decision fibre. This is the precise sense in which an envelope can
absorb the strongest available donor mechanisms while still identifying a missing
interface relation.

\paragraph{Corollary 1 (information-equivalent products must tie).}
Any donor product that reconstructs the transition tuple above and applies the
same decision rule is extensionally identical to $\Gamma$ on every context.

The ideal-product tie in Section~\ref{sec:revision} is therefore a required
negative boundary. It is evidence against an inherent centralization or
expressivity advantage, not an inconvenient null result.

\subsection{The optimal error of a lossy scientific interface}

Exact equivalence is a worst-case property. A distributional version
quantifies how much error remains when a donor interface merges contexts that
usually, but not always, require the same transition. Let $E\sim\mu$ on a
standard Borel context space, let $\phi$ be measurable into a standard Borel
interface space, let the output alphabet $\mathcal Y$ be finite,
and define the regular conditional probabilities
\[
p_y(z)=\Pr\{\Gamma(E)=y\mid\phi(E)=z\}.
\]

\paragraph{Theorem 2 (mixed-fibre Bayes lower bound).}
For every deterministic or randomized policy that uses only $\phi(E)$, its
zero--one transition error satisfies
\[
\Pr\{\widehat\Gamma\ne\Gamma(E)\}
\ \ge\
1-\mathbb E_{\phi(E)}\!\left[\max_{y\in\mathcal Y}p_y(\phi(E))\right].
\]
A fibrewise maximum-probability decoder attains equality. Consequently, zero
error is possible exactly when the conditional decision distribution is a
point mass almost surely.

\paragraph{Proof.}
Condition on $\phi(E)=z$. If a randomized decoder selects $y$ with probability
$q_y(z)$, its conditional success probability is
$\sum_y q_y(z)p_y(z)\leq\max_y p_y(z)$. Averaging and complementing gives the
bound. Concentrating $q$ on a maximizing label attains it. The final claim
follows because the expectation of the maximum is one exactly when the
conditional distribution is almost surely degenerate. \hfill$\square$

The right-hand side is the irreducible \emph{decision-fibre impurity} of the
interface. It supplies a population estimand for naturalistic evaluation:
the scientific question is not merely whether a composed system contains the
right named modules, but whether its visible interface collapses contexts with
different licensed transitions. This decision-theoretic reading is related
to Blackwell's comparison of experiments \cite{blackwell1953}; the contribution
here is its specialization to the typed, hard-gated transition object and the
empirical isolation of a particular omitted coupling, not the general idea of
comparing information structures.

\paragraph{Corollary 2 (component-essentiality criterion).}
Let the complete interface be $S=(S_1,\ldots,S_m)$ and let $S_{-j}$ omit one
component. Exact decisions are possible from $S_{-j}$ if and only if no two
contexts agree on $S_{-j}$ while requiring different $\Gamma$ decisions. If a
fibre contains two equiprobable decision classes, every policy omitting $S_j$
has conditional error at least $1/2$ on that fibre.

This turns every negative ablation into a new scientific question: characterize
the mixed fibres it creates, identify which missing relation splits them, and
test whether those fibres occur under a naturalistic acquisition distribution.
It does not turn a failed outcome into a positive result by relabeling it.

\paragraph{Corollary 3 (licence-coupling impossibility).}
Consider an interface that exposes candidate revisions, recovery results,
invariant checks, dependency impacts, and budgets, but omits $J_e$, the relation
coupling evidence to the layer authorized to change. No policy using only that
interface can agree with $\Gamma$ on every context. Under an equal mixture of
the two contexts below, every policy has error at least $1/2$.

\paragraph{Proof.}
Take $r_{\mathrm L}\prec r_{\mathrm H}$. Construct two contexts with identical
candidate, recovery, preservation, impact, and budget outputs. In $e_0$, the
evidence licenses only $r_{\mathrm L}$. In $e_1$, a registered exclusion
witness licenses only $r_{\mathrm H}$. Hence
$\Gamma(e_0)=r_{\mathrm L}$ and $\Gamma(e_1)=r_{\mathrm H}$, although the
reduced interface is identical. Theorem~2 gives the equal-mixture bound.
\hfill$\square$

The licence example is one instance of Theorems~1--2. The same analysis
applies to any omitted preservation, scope, measurement, authority, or resource
relation that changes the decision within an otherwise identical interface
fibre.

\subsection{Absorbing stronger theories without erasing old results}

A broad framework must be able to incorporate a stronger donor theory without
silently changing old conclusions. Let an old transition system on
$\mathcal R$ have admissible set $\mathcal A$ and front $\mathcal F$. Let an
extended system on $\mathcal R^+$ have admissible set $\mathcal A^+$ and front
$\mathcal F^+$, with an order embedding $i:\mathcal R\hookrightarrow
\mathcal R^+$ that preserves and reflects the old order and satisfies
$\mathcal A^+\cap i(\mathcal R)=i(\mathcal A)$. Write
$\mathcal N=\mathcal A^+\setminus i(\mathcal A)$ for the genuinely new
admissible transitions and assume the extended admissible order is well
founded.

\paragraph{Theorem 3 (conservative-envelope criterion).}
The extension preserves the old scientific front exactly,
$\mathcal F^+=i(\mathcal F)$, if and only if both conditions hold:
\begin{enumerate}
  \item no $n\in\mathcal N$ lies below an old minimum $i(f)$; and
  \item every $n\in\mathcal N$ lies above at least one old minimum $i(f)$,
  for some $f\in\mathcal F$.
\end{enumerate}

\paragraph{Proof.}
If the conditions hold, no old minimum gains a predecessor, while every new
admissible transition has an old admissible predecessor. The minimal elements
are therefore exactly the embedded old minima. Conversely, if the fronts are
equal, a new transition below an old minimum would destroy that minimum. And
because the extended order is well founded, every nonminimal new transition
has a minimal predecessor; equality of the fronts makes that predecessor an
embedded member of $\mathcal F$. \hfill$\square$

The criterion distinguishes assimilation from evasion. A donor mechanism that
only adds scientifically broader alternatives above an existing admissible
front can be absorbed conservatively. A donor that supplies a genuinely
narrower authorized transition must change the front; preserving the old claim
in that case would be scientifically wrong. Historical adverse results remain
attached to their original transition system, while a materially enlarged
system receives a new identity and a new test.

\paragraph{Proposition 2 (universal front representation).}
Fix any partial order $(\mathcal R,\preceq)$. Every antichain-valued map
$H:\mathcal E\to 2^{\mathcal R}$ is the minimal-front map of an epistemic
transition envelope.

\paragraph{Proof.}
For each context define the admissible set as the upward closure
$\mathcal A_H(e)=\{r:\exists h\in H(e),\ h\preceq r\}$. Every $h\in H(e)$ is
minimal because $H(e)$ is an antichain, and every other admissible element has
some $h$ strictly below it. Hence $\min\mathcal A_H(e)=H(e)$. \hfill$\square$

The proposition removes any dependence on the particular $K/W/M$ chain. The
same theory can represent incomparable transitions among questions,
representations, measurements, model classes, experimental designs, search
policies, and methods, provided their scientific change relation is made
explicit.

\subsection{What the present experiments establish}

The mutation-necessity experiment in Section~\ref{sec:necessity} tests one
five-predicate envelope and direct removals of its gates. The
revision-responsibility battery in Section~\ref{sec:revision} tests a finite set
of mixed fibres created by omitting the licence coupling and verifies exact
factorization with the information-equivalent product. The historical study
in Section~\ref{sec:hidden-formulation} remains the falsifier that exposed the
unsafe shortcut from diagnosis to high-level rewrite.

These studies do not estimate decision-fibre impurity in open science. The
wide successor therefore targets source-clustered naturalistic episodes across
eight causal families and four scientific domains, two source-disjoint waves,
and nine frozen comparator families. Until source snapshots, a strongest
runnable external comparator, changed semantic hosts, and independent scorer
custody are bound and the study is executed, that successor is a design rather
than evidence. The analytic theory is general; the present empirical authority
is not.
